\documentclass[12pt]{article}
\usepackage[utf8]{inputenc}
\usepackage[T2A]{fontenc}
\usepackage[russian]{babel}
\usepackage{amsmath,amssymb,amsthm}
\usepackage[a4paper,left=25mm,right=20mm,top=22mm,bottom=25mm]{geometry}
\usepackage{array}

\newtheorem{theorem}{Теорема}
\newtheorem{lemma}{Лемма}
\newtheorem{proposition}{Предложение}
\newtheorem{corollary}{Следствие}
\theoremstyle{definition}
\newtheorem{definition}{Определение}
\theoremstyle{remark}
\newtheorem{remark}{Замечание}

\newcommand{\R}{\mathbb{R}}
\newcommand{\Z}{\mathbb{Z}}
\newcommand{\E}{\mathbb{E}}
\newcommand{\diam}{\operatorname{diam}}
\newcommand{\dist}{\operatorname{dist}}

\title{Верхняя оценка $\chi(\R^4)\le 46$ для хроматического числа
четырёхмерного евклидова пространства%
\thanks{Черновик от 22.07.2026; аффилиация и данные автора подлежат уточнению.}}
\author{Леонид Иванов}
\date{22 июля 2026 г.}

\begin{document}
\maketitle

\begin{abstract}
Построена решёточная раскраска пространства $\R^4$ в $46$ цветов, при которой
одноцветные точки не реализуют ни одного расстояния из отрезка $[1,\ell]$ при
любом $\ell\le 1{,}004$. В частности, $\chi(\R^4)\le 46$: это улучшает верхнюю
оценку $\chi(\R^4)\le 49$, известную с 2006~г.
\cite{Ivanov2006,CoulsonPayne2007,ABPR}, и опровергает ожидание Армана,
Бондаренко, Прымака и Радченко (\cite{ABPR}, Open question~1) о том, что $49$
--- минимальный индекс решёточной (sublattice) раскраски $\E^4$. Раскраска
задаётся явной рациональной решёткой общего положения (ячейка Вороного ---
комбинаторный пермутоэдр с $30$ фасетами и $f$-вектором $(120,240,150,30)$),
найденной численной оптимизацией метрики по пространству квадратичных форм.
Неравенство $\ell<D(\Gamma)/\diam V_0$ доказано в \emph{точной рациональной
арифметике} (объемлющий многогранник из $30$ биссекторов, $120$ точных вершин,
опорные сертификаты) --- результат не зависит от вычислений с плавающей точкой.
Дополнительно приведена решётка, дающая $\chi\bigl(\R^4,[1,\ell]\bigr)\le48$
для более широкого интервала $\ell\le1{,}0396$, и вычислены точные нормированные
ширины запрещённых интервалов всех конструкций работы \cite{ABPR} в
размерностях $4\le n\le 9$.
\end{abstract}

\section{Введение}

Хроматическим числом $\chi(\R^n)$ называется наименьшее число цветов,
достаточное для такой раскраски точек $\R^n$, что любые две точки на расстоянии
$1$ окрашены различно (проблема Нельсона--Хадвигера; см. обзоры
\cite{Raigorodskii2000,Soifer}). Более общо, для отрезка $[1,\ell]$,
$\ell\ge 1$, через $\chi\bigl(\R^n,[1,\ell]\bigr)$ обозначим наименьшее число
цветов раскраски, при которой одноцветные точки не реализуют ни одного
расстояния из $[1,\ell]$ (хроматическое число с интервалом запрещённых
расстояний; систематическое изучение начато в \cite{Ivanov2011}). Очевидно
$\chi(\R^n)=\chi\bigl(\R^n,[1,1]\bigr)\le\chi\bigl(\R^n,[1,\ell]\bigr)$ при
$\ell\ge1$.

Известные границы в малых размерностях:
$5\le\chi(\R^2)\le 7$ \cite{deGrey,Soifer},
$6\le\chi(\R^3)\le 15$ \cite{Nechushtan,Coulson},
$9\le\chi(\R^4)$ \cite{ExooIsmailescu}. Верхняя оценка в размерности $4$
имеет следующую историю: $\chi(\R^4)\le 54$ (Радоичич, Тот \cite{RadoicicToth},
линейная раскраска на решётке $A_4^*$); $\chi(\R^4)\le 49$ --- анонсировано в
заметке \cite{Ivanov2006} и упомянуто в \cite{CoulsonPayne2007}; полное
доказательство, вместе с исчерпывающей проверкой того, что $49$ ---
минимальный индекс допустимой подрешётки для решётки $D_4$ (а $54$ --- для
$A_4^*$), дано Арманом, Бондаренко, Прымаком и Радченко \cite{ABPR}. Там же
(Open question~1) авторы сформулировали ожидание, что улучшение невозможно:
\emph{<<We believe that the answer is negative for $n=4,5$, i.e.
$\min\chi_s(\E^4,\Lambda,\Lambda',\ell)=49$~\ldots, where the minima are taken
over all full rank lattices $\Lambda$, sublattices $\Lambda'$ and all
$\ell>0$>>}.

Основной результат настоящей работы опровергает это ожидание.

\begin{theorem}\label{thm:main}
Пусть $\ell_0=1{,}004$. Тогда
\[
\chi\bigl(\R^4,[1,\ell]\bigr)\;\le\;46
\qquad\text{при всех }1\le\ell\le\ell_0 .
\]
В частности, $\chi(\R^4)\le 46$.
\end{theorem}

Раскраска в теореме~\ref{thm:main} --- решёточная (классы смежности подрешётки
индекса $46$ раскрашивают ячейки Вороного явно заданной рациональной решётки
$\Lambda_\star\subset\R^4$), т.е. в терминологии \cite{ABPR}
$\min_{\Lambda,\Lambda',\ell}\chi_s(\E^4,\Lambda,\Lambda',\ell)\le 46$.
Существенно, что $\Lambda_\star$ --- решётка \emph{общего положения}
(генерическая): у неё лишь одна пара минимальных векторов, а ячейка Вороного
имеет максимально возможные $30$ фасет. Все ранее рассматривавшиеся
конструкции в $\R^4$ опирались на классические решётки ($D_4$, $A_4^*$,
$\Z^4$); исчерпывающие переборы работы \cite{ABPR} покрывали только $D_4$ и
$A_4^*$ и потому не могли обнаружить $\Lambda_\star$.

Метод даёт целое семейство генерических решёток, последовательно улучшающих
границу: индексы $48$, $47$ (следствие) и $46$ достигаются, тогда как наилучшая
найденная нормированная ширина при индексе $45$ равна лишь $0{,}9955<1$
(раздел~\ref{sec:frontier}). Нижняя оценка Кулсона для решёточных раскрасок
$\chi_s(\E^n,\Lambda,\Lambda',\ell)\ge 2^{n+1}-1$ \cite{Coulson,ABPR} в
размерности $4$ равна $31$ и не препятствует значению $46$.

\section{Метод решёточных раскрасок}\label{sec:method}

Напомним конструкцию (Л.\,Л.~Иванов \cite{Ivanov2011}; в эквивалентной форме
--- \cite[предложение~5]{ABPR}). Пусть $\Lambda\subset\R^n$ --- решётка полного
ранга с ячейкой Вороного
$V_0=\{x:\ \|x\|\le\|x-v\|\ \forall v\in\Lambda\}$, и пусть
$\Gamma\subseteq\Lambda$ --- подрешётка индекса $k=[\Lambda:\Gamma]$. Раскрасим
ячейку $V_\Lambda(v)$, $v\in\Lambda$, цветом класса смежности $v+\Gamma$ ---
получится $k$-цветная периодическая раскраска $\R^n$.

\begin{lemma}[Вороной; см. \cite{Voronoi}]\label{lem:sym}
Ячейка $V_0$ --- выпуклый центрально-симметричный многогранник, и все её грани
центрально-симметричны.
\end{lemma}

\begin{lemma}[лемма о середине; {\cite[лемма~2]{Ivanov2011}}]\label{lem:mid}
Минимальное расстояние между ячейками $V_\Lambda(a)$ и $V_\Lambda(b)$
реализуется на паре точек, симметричных относительно середины отрезка $[a,b]$;
как следствие, для $v\in\Lambda$
\[
D(v):=\dist\bigl(V_0,\,v+V_0\bigr)=2\,\dist\bigl(v/2,\,V_0\bigr).
\]
\end{lemma}

Положим
\[
D(\Gamma)=\min_{v\in\Gamma\setminus\{0\}}D(v),\qquad
d(\Lambda,\Gamma)=\frac{D(\Gamma)}{\diam V_0}.
\]
Две точки одного цвета лежат либо в одной ячейке (расстояние $\le\diam V_0$),
либо в ячейках $V_\Lambda(a)$, $V_\Lambda(b)$ с $a-b\in\Gamma\setminus\{0\}$
(расстояние $\ge D(\Gamma)$). Поэтому множество реализуемых одноцветных
расстояний не пересекает \emph{открытый} интервал $\bigl(\diam V_0,
D(\Gamma)\bigr)$.

\begin{proposition}[{ср. \cite[Prop.~5]{ABPR}}]\label{prop:crit}
Если $d(\Lambda,\Gamma)>1$, то $\chi\bigl(\R^n,[1,\ell]\bigr)\le[\Lambda:\Gamma]$
при всех $\ell<d(\Lambda,\Gamma)$; в частности $\chi(\R^n)\le[\Lambda:\Gamma]$.
\end{proposition}

\begin{proof}
Выберем масштаб $s>0$ так, что $s\,\diam V_0<1$ и $s\,D(\Gamma)>\ell$; это
возможно тогда и только тогда, когда $\ell<D(\Gamma)/\diam V_0=d(\Lambda,\Gamma)$.
Тогда для решётки $s\Lambda$ отрезок $[1,\ell]$ целиком лежит в интервале
$\bigl(s\diam V_0,\,sD(\Gamma)\bigr)$, свободном от одноцветных расстояний,
--- раскраска запрещает весь $[1,\ell]$. Неравенство для $\chi(\R^n)$
получается при $\ell=1$.
\end{proof}

Практическое вычисление $D(\Gamma)$ опирается на неравенство
$D(v)\ge\|v\|-\diam V_0$ (ячейка лежит в шаре радиуса $\tfrac12\diam V_0$):
достаточно перебрать конечное множество векторов
$\{v\in\Gamma:\ \|v\|<D_{\text{тек}}+\diam V_0\}$, где $D_{\text{тек}}$ ---
текущая верхняя оценка минимума; перебор полон при любом выборе базиса.

\section{Конструкция}\label{sec:construction}

Определим решётку $\Lambda_\star=B^{\mathsf T}\Z^4$, где $B$ --- фактор
Холецкого ($Q=BB^{\mathsf T}$, строки $B$ --- базисные векторы) следующей
рациональной матрицы Грама:
\begin{equation}\label{eq:gram}
Q=\begin{pmatrix}
\tfrac{176}{15} & \tfrac{19}{6}  & -\tfrac{79}{10} & -\tfrac{143}{24}\\[2pt]
\tfrac{19}{6}   & \tfrac{25}{2}  & -\tfrac{19}{4}  & -\tfrac{97}{17}\\[2pt]
-\tfrac{79}{10} & -\tfrac{19}{4} & 24              & -\tfrac{89}{21}\\[2pt]
-\tfrac{143}{24}& -\tfrac{97}{17}& -\tfrac{89}{21} & \tfrac{53}{3}
\end{pmatrix},
\qquad
14280\,Q\in\mathrm{Mat}_4(\Z).
\end{equation}
Подрешётку $\Gamma_\star\subset\Lambda_\star$ индекса $46$ зададим матрицей
перехода в эрмитовой нормальной форме (строки коэффициентов относительно
базиса $B$):
\begin{equation}\label{eq:hnf}
T=\begin{pmatrix}
1&0&0&44\\
0&1&0&32\\
0&0&1&42\\
0&0&0&46
\end{pmatrix},
\qquad \det T=46 .
\end{equation}

Числовые инварианты (величины $\lambda_1^2$, $D^2$, $\diam^2 V_0$ рациональны,
поскольку рациональна $Q$; см. точные значения в разделе~\ref{sec:comp}):
\[
\lambda_1(\Lambda_\star)^2=\tfrac{176}{15},\qquad
\diam^2 V_0\approx 32{,}8519,\qquad
D(\Gamma_\star)^2\approx 33{,}1858,
\]
\begin{equation}\label{eq:d}
d(\Lambda_\star,\Gamma_\star)=\frac{D(\Gamma_\star)}{\diam V_0}
=1{,}005068981\ldots>\ell_0=1{,}004 .
\end{equation}
Ячейка Вороного $V_0(\Lambda_\star)$ имеет $30$ фасет (15 пар релевантных
векторов --- максимум в размерности~4) и $f$-вектор $(120,240,150,30)$
комбинаторного пермутоэдра; у $\Lambda_\star$ единственная пара минимальных
векторов, отношение покрытие/упаковка $2R/\lambda_1\approx1{,}673$. Решётка не
изометрична ни одной из классических ($D_4$, $A_4^*$, $\Z^4$, $A_4$).

Теорема~\ref{thm:main} следует из предложения~\ref{prop:crit} и неравенства
\eqref{eq:d} (сертификат которого --- в разделе~\ref{sec:comp}).

\begin{remark}[численный оптимум]\label{rem:float}
Матрица \eqref{eq:gram} --- рационализация (со знаменателями $\le24$)
численного оптимума задачи $\max_Q d(Q,46)$, найденного методом Нелдера--Мида.
Сам оптимум даёт несколько бо́льшую ширину, $d(46)=1{,}0054054\ldots$ (значение
воспроизведено двумя независимыми программами), причём оптимальная подрешётка
задаётся \emph{той же} матрицей перехода \eqref{eq:hnf}. Для строгой
формулировки теоремы~\ref{thm:main} используется рациональная форма
\eqref{eq:gram}, допускающая точный арифметический сертификат.
\end{remark}

\section{Точный сертификат неравенства \eqref{eq:d}}\label{sec:comp}

Помимо согласия трёх независимых численных алгоритмов вычисления
$\dist(v/2,V_0)$ (перебор по вершинам ячейки методом Гилберта--Джонсона--Кирти
с апостериорным контролем зазора двойственности; каскад ортогональных проекций
по граням; решение задачи квадратичного программирования по полупространствам)
--- все дают $d=1{,}005068981$ с совпадением в девяти значащих цифрах ---
неравенство $d\ge\ell_0$ проверено в \emph{точной рациональной арифметике}.

Пусть $P$ --- пересечение $30$ полупространств-биссекторов
$\{x:\ \langle x,v\rangle_Q\le\tfrac12\langle v,v\rangle_Q\}$ по $15$ парам
целочисленных векторов $\pm v$ (кандидатов в релевантные, найденных
CVP-перебором классов смежности $\Lambda_\star/2\Lambda_\star$). Каждый
биссектор содержит $V_0$, поэтому $V_0\subseteq P$; следовательно
$\diam V_0\le\diam P$ и $\dist(p,V_0)\ge\dist(p,P)$ для любой точки~$p$.
\begin{enumerate}
\item[(а)] \emph{Верхняя оценка диаметра.} Все вершины $P$ перечислены точно
(перебор $\binom{30}{4}=27\,405$ рациональных систем $4\times4$ с проверкой
всех неравенств; вершин ровно $120$), и
\[
\diam^2 V_0\le\diam^2 P=\frac{20390598460067808}{620682234555961}
=32{,}851912500\ldots
\]
\item[(б)] \emph{Нижняя оценка $D(\Gamma_\star)$.} Для каждого
$v\in\Gamma_\star$ в доказуемо полном окне кандидатов
($\|v\|<D_{\text{тек}}+\diam V_0$): либо точно проверяется
$\langle v,v\rangle_Q\ge W^2$, где рациональное $W$ удовлетворяет
$W^2\ge(1+\ell_0)^2\diam^2 P$ (тогда $D(v)\ge\|v\|-\diam V_0\ge\ell_0\diam V_0$
автоматически), либо --- для $11$ коротких кандидатов --- применяется опорное
неравенство
\[
\dist(v/2,P)\ \ge\ \frac{\langle v/2,u\rangle_Q-\max_{w\in\mathrm{vert}\,P}
\langle w,u\rangle_Q}{\|u\|_Q}
\]
с рациональным направлением $u$ (округление оптимального направления
квадратичной задачи проекции). Все величины --- точные дроби; в итоге
\[
D(\Gamma_\star)^2\ \ge\ \frac{496626050089}{14965013040}
=33{,}185808041\ldots
\]
\item[(в)] \emph{Вердикт.} $\ell_0^2\,\diam^2 P=33{,}115253430\ldots$, и
поскольку $33{,}185808041>33{,}115253430$, получаем
$D(\Gamma_\star)^2\ge\ell_0^2\diam^2 P\ge\ell_0^2\diam^2 V_0$, т.е.
$d(\Lambda_\star,\Gamma_\star)\ge\ell_0$.
\end{enumerate}
Сертификат не использует плавающую арифметику нигде, кроме \emph{выбора}
кандидатов и направлений (на строгость не влияющего). Дополнительно выполнен
исчерпывающий перебор всех $190\,800$ подрешёток индекса $46$ решётки
$\Lambda_\star$, подтверждающий, что \eqref{eq:hnf} доставляет максимум $d$ на
этом индексе.

\paragraph{Как найдена решётка.} Максимизировалась функция
$t\mapsto d\bigl(Q(t),k\bigr)$ методом Нелдера--Мида по пространству
положительно определённых форм (параметризация Холецкого): сначала для $k=48$
--- со стартом во вторичном конусе генерического типа Делоне (в обозначениях
\cite{VWZ} --- тип $111{-}$; у его стандартного представителя $d(48)=0{,}99941$),
затем последовательным спуском по $k$ с многократными рестартами вблизи
найденных оптимумов. Скрипты поиска, верификации и точного сертификата, а также
все промежуточные данные доступны у автора (пакеты \texttt{voronoi4d} и
\texttt{combigeo}, версия~1.1.0, с многопоточным перебором подрешёток).

\section{Более широкий интервал: $\chi(\R^4,[1,\ell])\le48$}\label{sec:wide}

Та же оптимизация при индексе $48$ даёт заметно бо́льшую ширину запрещённого
интервала. Для рациональной решётки с матрицей Грама
\[
Q_{48}=\begin{pmatrix}
\tfrac{137}{12} & \tfrac{13}{7}  & -\tfrac{61}{10} & -\tfrac{18}{5}\\[2pt]
\tfrac{13}{7}   & 7              & -\tfrac{25}{9}  & -\tfrac{37}{12}\\[2pt]
-\tfrac{61}{10} & -\tfrac{25}{9} & 12              & -\tfrac{17}{11}\\[2pt]
-\tfrac{18}{5}  & -\tfrac{37}{12}& -\tfrac{17}{11} & \tfrac{44}{5}
\end{pmatrix}
\]
и подрешётки с матрицей перехода
$\mathrm{diag\text{-}HNF}(1,1,1,48)$ с наддиагональю $(28,16,30)$ в последнем
столбце получается $d=1{,}039657681\ldots$; тем же точным сертификатом
(здесь $\lambda_1^2=7$, $D^2\ge20{,}689971$, $\diam^2 P=19{,}141641$)
доказывается
\[
\chi\bigl(\R^4,[1,\ell]\bigr)\le48\qquad\text{при }\ell\le1{,}0396 .
\]
Численный оптимум при индексе $48$ достигает $d=1{,}0432696\ldots$
Таким образом, при небольшом увеличении числа цветов ($46\to48$) существенно
расширяется запрещённый интервал ($1{,}004\to1{,}0396$) --- компромисс, не
отражённый в предыдущей литературе.

\section{Граница метода и сопутствующие результаты}\label{sec:frontier}

\paragraph{Насколько мал индекс?} Оптимизация формы при индексе $45$ даёт
наилучшую найденную ширину $d(45)=0{,}9955<1$ (при $\ge10^3$ рестартах на
пуле процессов), а при индексах $44$, $43$ --- ещё меньшие значения. Хотя это
не исключает существования подходящей решётки, мы склонны предполагать, что
$46$ близко к пределу решёточных раскрасок $\R^4$; строгое доказательство
минимальности остаётся открытым (ср. постановку Open question~1 в \cite{ABPR}).

\paragraph{Точные ширины конструкций \cite{ABPR} в размерностях $4\le n\le9$.}
Тем же инструментарием вычислены нормированные ширины запрещённых интервалов
всех явных конструкций работы \cite{ABPR} (в самой работе доказывалось только
$d\ge1$); каждая строка даёт оценку $\chi(\R^n,[1,\ell])\le k$ при $\ell<d$:
\begin{center}
\begin{tabular}{c c l l}
\hline
$n$ & $k$ & конструкция & $d$ (точно)\\
\hline
4 & 49  & $(3+\omega)D_4$ & $\sqrt{7/6}=1{,}080123\ldots$\\
4 & 54  & $A_4^*$ & $\sqrt{23/20}=1{,}072381\ldots$\\
5 & 140 & $A_5^*$, матрица $C_5$ из \cite{ABPR} & $\sqrt{39/35}=1{,}055597\ldots$\\
6 & 343 & $(3+\omega)E_6^*$ & $\sqrt{7/6}$\\
7 & 1372& $E_7^*$, матрица $C_7$ из \cite{ABPR} & $\sqrt{17/14}=1{,}101946\ldots$\\
8 & 2401& $(3+\omega)E_8$ & $\sqrt{7/6}$\\
9 & 17253& $A_9^*$, матрица $C_9$ из \cite{ABPR} & $13/\sqrt{165}=1{,}012048\ldots$\\
\hline
\end{tabular}
\end{center}

Для эйзенштейновской конструкции $\Gamma=(3+\omega)\Lambda$ ($|3+\omega|^2=7$,
в решётках с $\Z[\omega]$-структурой) во всех проверенных случаях
($n=2,4,6,8$) выполняется точное равенство
$D\bigl((3+\omega)\Lambda\bigr)^2=\tfrac73\lambda_1(\Lambda)^2$, т.е.
$d=\sqrt{7/3}\,/\,(2R/\lambda_1)$. При $2R/\lambda_1=\sqrt2$ (решётки $D_4$,
$E_6^*$, $E_8$, $\Lambda_{24}$) отсюда всегда $d=\sqrt{7/6}$; в частности, в
предположении сохранения равенства при $n=24$ получаем гипотетическую оценку
$\chi\bigl(\R^{24},[1,\sqrt{7/6}]\bigr)\le 7^{12}$. Для решётки Кокстера--Тодда
$K_{12}$ ($2R/\lambda_1=\sqrt{8/3}$) та же формула даёт $d\approx0{,}935<1$:
конструкция теоремы~3 из \cite{ABPR} на $K_{12}$ не проходит (частичный
отрицательный ответ на их Open question~2). Исчерпывающим перебором всех
$1\,424\,115\,231$ подрешёток индекса $140$ решётки $A_5^*$ установлено, что
конструкция $C_5$ из \cite{ABPR} оптимальна на этом индексе и по ширине:
$\max d=\sqrt{39/35}$.

\paragraph{Улучшения в $\R^3$.} Оптимизация метрики по всему шестимерному
пространству форм улучшает таблицу работы \cite{Ivanov2011} (каждое значение
--- явная решётка и подрешётка):
\[
\begin{array}{c|ccccccccccc}
k & 16 & 20 & 21 & 22 & 23 & 24 & 26 & 28 & 29 & 30 & 31\\
\hline
d & 1{,}0254 & 1{,}1375 & 1{,}1801 & 1{,}1893 & 1{,}2140 & 1{,}3158 &
1{,}2637 & 1{,}4596 & 1{,}3757 & 1{,}5253 & 1{,}4288
\end{array}
\]
(в \cite{Ivanov2011}: $k=21\mapsto1{,}133$, $k=23\mapsto1{,}137$,
$k=24\mapsto1{,}303$; значения при $k=16$, $20$, $22$, $26$, $28$--$31$ ранее
не публиковались).

\section{Открытые вопросы}

\begin{enumerate}
\item Является ли $46$ минимальным числом цветов решёточной раскраски $\R^4$?
Наши эксперименты ($\max d\approx0{,}9955$ при $k=45$) позволяют предположить
утвердительный ответ, но доказательства нет.
\item Найти аналитическое описание оптимальных генерических решёток для
индексов $46$ и $48$ (предельных точек максимизации~$d$).
\item Останется ли значение $140$ минимальным в $\R^5$, а $343$ --- в $\R^6$,
при оптимизации по генерическим решёткам? (Наш рандомизированный поиск в
$\R^6$ по методике \cite{ABPR} улучшений не дал.)
\item Доказать равенство $D((3+\omega)\Lambda)^2=\tfrac73\lambda_1^2$ для всех
эйзенштейновских решёток или описать область его справедливости.
\end{enumerate}

\subsection*{Благодарности}
Вычисления выполнены пакетами \texttt{voronoi4d} и \texttt{combigeo} (версия
1.1.0) автора. Автор признателен авторам работы \cite{ABPR} за ясную
постановку открытых вопросов.

\begin{thebibliography}{20}

\bibitem{ABPR}
A.~Arman, A.~Bondarenko, A.~Prymak, D.~Radchenko,
\emph{Upper bounds on chromatic number of $\E^n$ in low dimensions},
Electron. J. Combin. \textbf{31} (2024), no.~2, \#P2.35;
arXiv:2112.13438.

\bibitem{CoulsonPayne2007}
D.~Coulson, M.~S.~Payne,
\emph{A dense distance 1 excluding set in $\R^3$},
Austral. Math. Soc. Gaz. \textbf{34} (2007), no.~2, 97--102.

\bibitem{Coulson}
D.~Coulson,
\emph{A 15-colouring of 3-space omitting distance one},
Discrete Math. \textbf{256} (2002), 83--90.

\bibitem{deGrey}
A.~D.~N.~J. de~Grey,
\emph{The chromatic number of the plane is at least~5},
Geombinatorics \textbf{28} (2018), no.~1, 18--31.

\bibitem{DSMMV}
M.~Dutour~Sikiri\'c, D.~Madore, P.~Moustrou, F.~Vallentin,
\emph{Coloring the Voronoi tessellation of lattices},
J. London Math. Soc. (2) \textbf{104} (2021), 1135--1171.

\bibitem{ExooIsmailescu}
G.~Exoo, D.~Ismailescu,
\emph{On the chromatic number of $\R^4$},
Discrete Comput. Geom. \textbf{52} (2014), 416--423.

\bibitem{Ivanov2006}
Л.~Л.~Иванов,
\emph{Оценка хроматического числа пространства $\R^4$},
УМН \textbf{61} (2006), №~5, 181--182.

\bibitem{Ivanov2011}
Л.~Л.~Иванов,
\emph{О хроматических числах $\R^2$ и $\R^3$ с интервалами запрещённых
расстояний},
Вестник РУДН, сер. Математика. Информатика. Физика (2011), №~1, 14--29.

\bibitem{LarmanRogers}
D.~G.~Larman, C.~A.~Rogers,
\emph{The realization of distances within sets in Euclidean space},
Mathematika \textbf{19} (1972), 1--24.

\bibitem{Nechushtan}
O.~Nechushtan,
\emph{On the space chromatic number},
Discrete Math. \textbf{256} (2002), 499--507.

\bibitem{RadoicicToth}
R.~Radoi\v{c}i\'c, G.~T\'oth,
\emph{Note on the chromatic number of the space},
Discrete and Computational Geometry, Algorithms Combin., vol.~25,
Springer, Berlin, 2003, 695--698.

\bibitem{Raigorodskii2000}
А.~М.~Райгородский,
\emph{О хроматическом числе пространства},
УМН \textbf{55} (2000), №~2, 147--148.

\bibitem{Soifer}
A.~Soifer,
\emph{The Mathematical Coloring Book},
Springer, New York, 2009.

\bibitem{Voronoi}
Г.~Ф.~Вороной,
\emph{Собрание сочинений}, т.~2,
Изд-во АН УССР, Киев, 1952.

\bibitem{VWZ}
F.~Vallentin, S.~Wei\ss bach, M.~C.~Zimmermann,
\emph{The chromatic number of 4-dimensional lattices},
arXiv:2407.03513 (2024).

\end{thebibliography}

\end{document}
